% -----------------------------------------------------------------------------
% This file is automatically generated.
% Do not edit manually.
% Source: notebooks/support/regression-analysis/support.Rmd
% -----------------------------------------------------------------------------


\noindent To assess multicollinearity for a model without interaction terms, such as \ttblue{mod.1}, we use the standard \ttblue{vif()} function.

\par\addvspace{\topsep}
\begingroup
\fvset{listparameters={%
  \setlength{\topsep}{0pt}%
  \setlength{\partopsep}{0pt}%
  \setlength{\parsep}{0pt}%
  \setlength{\itemsep}{0pt}%
}}
\begin{Verbatim}[breaklines=true,formatcom=\color{ada_blue}]
vif(ussteamco.mod.1)
\end{Verbatim}
\begin{Verbatim}[breaklines=true,formatcom=\color{ada_light_blue}]
production     coolDD     heatDD
  2.664278   5.599833   3.732744
\end{Verbatim}
\endgroup
\par\addvspace{\topsep}

\noindent A common rule of thumb is that VIF values greater than 10 indicate substantial multicollinearity. Since none of the VIF values exceeds this threshold, we conclude that there is no evidence of serious variance inflation in \ttblue{mod.1}.

\noindent For models that include interaction terms, such as \ttblue{mod.3}, we use the \ttblue{vif()} function with the additional option \ttblue{type = "predictor"}. This reports the Generalized Variance Inflation Factor (GVIF) and its adjusted value, \(GVIF^{1/(2df)}\).
